\chapter{Especificaciones tecnicas}
\thispagestyle{memoria}\pagestyle{memoria}

\section{Disposiciones generales}
Los materiales seran nuevos, normalizados, compatibles entre si y adecuados para
la tension, corriente, temperatura, altitud, humedad, radiacion solar y presencia
de hidrocarburos del lugar. Cada suministro tendra ficha, marcado, certificado y
manual cuando la criticidad lo requiera. La aceptacion de una marca no elimina la
obligacion de cumplir el CNE-U, la EM.010 y las reglas sectoriales aplicables.

Las dimensiones, potencias y modelos declarados como criterios de diseño deben
confirmarse antes de compra. En caso de discrepancia entre plano, cuadro y ficha,
se detendra el montaje de la partida afectada y se emitira una consulta tecnica;
no se resolvera reduciendo proteccion, seccion o grado de seguridad.

\section{Conductores y canalizaciones}
Los conductores seran de cobre, 0.6/1 kV donde corresponda, aislamiento LSZH o resistente a hidrocarburos, temperatura y humedad segun el ambiente. No se reducira el neutro del alimentador general. El PE sera continuo, verde/amarillo, sin fusibles ni interruptores. Las canalizaciones enterradas seran pesadas, selladas contra ingreso de combustible y agua, con registros accesibles fuera de zonas clasificadas siempre que sea posible.

En areas peligrosas solo se instalaran equipos, cajas, prensaestopas, sellos y sistemas de cableado certificados para la zona, grupo de gas y clase de temperatura aplicables. Ninguna marca de referencia sustituye la certificacion del equipo realmente adquirido.

Los conductores se tenderan continuos entre cajas o terminales, sin empalmes en
tuberias. Se dejara reserva ordenada en tableros y equipos, se respetara el radio
de curvatura y se identificaran fases R--S--T, neutro N y proteccion PE en ambos
extremos. La ocupacion de ductos, agrupamiento y correccion termica se verificaran
con el cable finalmente adquirido.

Las zanjas electricas se replantearan antes de excavar. Los ductos tendran cama,
proteccion mecanica, pendiente o drenaje cuando corresponda, cinta de advertencia
y registros que permitan mantenimiento sin abrir la zona de dispensacion. Los
cruces con tuberias de combustible se resolveran con separacion y detalle aprobado
por las especialidades responsables.

\section{Cajas, sellos y uniones}
Las cajas interiores comunes seran metalicas o no propagadoras de llama, con tapa
y volumen suficiente. Las exteriores tendran grado de proteccion compatible con
polvo, lluvia y lavado. En el limite de un area clasificada se instalaran sellos
certificados en la posicion exigida por el sistema de cableado y la ficha del
equipo. No se admiten adaptadores improvisados, roscas dañadas, entradas abiertas
ni prensaestopas sin certificacion.

\section{Tableros y protecciones}
TGE, TDF, TDE y tableros administrativos seran metalicos, con barras separadas N/PE, rotulado permanente, reserva fisica y diagrama. El interruptor principal sera 4P-80 A; los alimentadores se protegen segun IE-05. Los circuitos derivados usan RCBO de hasta 30 mA para reducir disparos comunes y mejorar continuidad.

El poder de corte de 25 kA en cabecera y 10 kA en derivados es un minimo conservador del anteproyecto. Debe recalcularse con la Icc de Electro Puno y verificarse selectividad, filiacion y energia pasante con una sola familia coordinada.

Cada tablero tendra esquema unifilar plastificado, directorio de circuitos,
borneras, barrera contra contactos, señal de riesgo y al menos 20\% de reserva
fisica. La barra de neutro permanecera aislada de la envolvente y separada de la
barra PE. Los tornillos se ajustaran al par indicado por fabricante y se
registrara la inspeccion termografica o de calentamiento durante puesta en marcha.

\section{Motores y cargas especiales}
Las STP, el compresor y las bombas de servicio tendran seccionamiento identificable,
proteccion contra cortocircuito, sobrecarga y falla a tierra, y control compatible
con la corriente de arranque. El ajuste se hara con la corriente de placa, no con
la potencia adoptada del anteproyecto. La perdida de una fase, el rearranque no
intencional y la actuacion del paro de emergencia se comprobaran funcionalmente.

\section{Bombas, surtidores y paro de emergencia}
Las cuatro STP se consideran de 1.5 hp, 208--230 V, hasta 11 A, como familia de referencia. Cada motor tendra proteccion contra cortocircuito, sobrecarga y falla a tierra coordinada con su fabricante. Los seis cabezales de surtidor se alimentan a 220 V mediante UPS senoidal, sin asumir que la electronica de referencia define el equipo final.

Los dispositivos de paro remoto desenergizaran bombas y surtidores sin anular alarma, comunicaciones de emergencia ni alumbrado seguro. Su ubicacion debe ser visible, accesible y fuera de la envolvente peligrosa; los controles ordinarios y avisos se coordinan a mas de 3 m de llenado y venteos conforme al criterio sectorial documentado.

El cableado de datos, pulsos y control del sistema de combustible se separara de
potencia y de fuentes de interferencia. Las barreras, interfaces y UPS se
seleccionaran con el fabricante del surtidor y del ATG. Toda entrada a la zona de
combustible conservara el metodo de proteccion certificado del conjunto.

\section{Alumbrado y dispositivos de utilizacion}
Las luminarias de marquesina y exterior tendran envolvente para intemperie,
resistencia mecanica, factor de potencia y proteccion contra sobretensiones
adecuados. Cuando una luminaria quede dentro de la zona clasificada, debera portar
marcado compatible; no basta un grado IP alto. Las luminarias interiores y de
emergencia se montaran de forma mantenible y se probaran con perdida simulada de
red.

Los tomacorrientes seran 2P+T con contacto de proteccion efectivo. En exteriores
se usaran tapas y envolventes para intemperie. No se ubicaran tomacorrientes
ordinarios dentro de areas peligrosas. Todos los circuitos de utilizacion tendran
proteccion diferencial de hasta 30 mA y rotulo de origen.

\section{Puesta a tierra, rayo y sobretensiones}
Se propone anillo de cobre desnudo de 35 mm\(^2\), ocho electrodos como punto de partida y union equipotencial de tanques, tuberias, surtidores, marquesina, tableros, grupo y todas las masas. El limite del CNE-U 060-712 es 25 ohm; se adopta objetivo de 10 ohm, sujeto a resistividad y medicion con telurimetro.

La proteccion contra el rayo se plantea como sistema convencional coordinado y no como electrodo aislado independiente de la instalacion. El riesgo, numero de captadores y separacion se verificaran con IEC 62305. Se instalaran SPD Tipo 1+2 en TGE y Tipo 2 aguas abajo, coordinados con el sistema de puesta a tierra.

Las conexiones enterradas del anillo se ejecutaran con conectores certificados o
soldadura exotermica segun el detalle aprobado. Las uniones visibles quedaran
accesibles para inspeccion. Los tanques, tuberias, islas, surtidores, marquesina,
grupo y tableros se enlazaran sin crear electrodos aislados. La resistencia medida
se registrara con fecha, metodo, instrumento y condiciones del suelo.

\section{Grupo electrogeno, ATS y UPS}
El grupo sera trifasico, 380/220 V, 60 Hz, potencia standby no menor de la
seleccion calculada y con confirmacion escrita del derrateo a la altitud del sitio.
El ATS sera de cuatro polos para conmutar el neutro y tendra enclavamiento que
impida paralelismo no autorizado. Se verificaran secuencia, tiempos, retorno de
red, parada y capacidad de arranque secuencial.

Las UPS seran online o de desempeño equivalente, onda senoidal, bypass y alarma.
La UPS-FUEL atendera exclusivamente electronica/control de combustible y la
UPS-IT atendera POS, comunicaciones y CCTV. La autonomia minima de 15 minutos se
demostrara con la carga real y baterias nuevas.

\section{Identificacion y documentacion}
Cada circuito, tablero, cable, boton, electrodo y equipo tendra codigo coincidente
con IE-01 a IE-06. Al finalizar se entregaran planos conforme a obra, cuadros
actualizados, fichas, certificados, ajustes de protecciones y protocolos. No se
retiraran advertencias de anteproyecto hasta que exista revision profesional y
documentacion de campo suficiente.

\section{Pruebas minimas}
Antes de energizar se mediran continuidad del PE y enlaces equipotenciales, resistencia de aislamiento, polaridad, resistencia de puesta a tierra, impedancia de lazo, secuencia de fases, operacion de RCBO, selectividad funcional, ATS, paro de emergencia, UPS, grupo, niveles de iluminacion y funcionamiento de cada equipo. Los resultados se registraran en protocolos firmados por personal competente.
